% Verbatim from /Users/alingalatan/Downloads/cubic_chaos_partial_results (2).zip:six_dimensional_proof.tex (16 November 2026); GPT (a ChatGPT session of the user), relayed by the user.
% Intended filepath: six_dimensional_proof.tex
\documentclass[11pt]{article}
\usepackage[margin=1in]{geometry}
\usepackage{amsmath,amssymb,amsthm}
\newtheorem{theorem}{Theorem}
\newcommand{\E}{\mathbb E}
\title{A strict Gaussian $L^1/L^2$ bound for cubic chaos in at most six dimensions}
\author{}
\date{}
\begin{document}
\maketitle
\noindent\textbf{Status.} The proof is rigorous and uses only elementary
angular inequalities, Gaussian polar decomposition, and the supplied exact
rational certificate. It does not settle arbitrary dimension.

\begin{theorem}
Every real cubic Wiener-chaos polynomial depending on at most six independent
standard real Gaussian coordinates satisfies
\[
 \boxed{\E|h|<\frac{83\sqrt5}{75\pi}\|h\|_2
       <\sqrt{\frac2\pi}\|h\|_2\qquad(h\ne0).}
\]
\end{theorem}

\begin{proof}
Embed lower-dimensional polynomials into six dimensions by adding unused
coordinates. Write a standard Gaussian vector as $R\Omega$, where
$\Omega$ is uniform on $S^5$, $R\sim\chi_6$, and they are independent.
Every homogeneous cubic has a unique decomposition into a harmonic cubic
$H_3$ plus $r^2 H_1$, with $H_1$ linear. Wick subtraction therefore gives
\[
 h(R\Omega)=R^3H_3(\Omega)+R(R^2-8)H_1(\Omega),
\]
since $\Delta(r^2H_1)=16H_1$ in dimension six. The two spherical harmonics
are orthogonal on $S^5$.

For every odd polynomial of degrees one and three on a sphere,
\begin{equation}\label{eq:angular}
 \E_\Omega|P(\Omega)|\leq\beta(\E_\Omega P(\Omega)^2)^{1/2},
 \qquad \beta=\frac{\sqrt{80}}{3\pi}.
\end{equation}
Indeed restrict $P$ to a random Haar-distributed great circle. On each such
circle it has frequencies only $\pm1,\pm3$, so the angular lemma proved in
\texttt{binary\_cubic\_proof.tex} applies. Average the circle inequality and
apply Cauchy--Schwarz over the random circle. A uniform point on a Haar
great circle is uniform on the sphere, proving \eqref{eq:angular}.

Normalize $\|h\|_2=1$ and let $\lambda\in[0,1]$ be the share of its variance
from the $R^3H_3$ term. Put $T=R^2/2\sim\operatorname{Gamma}(3,1)$, whose
density is $t^2e^{-t}/2$. Spherical orthogonality and elementary gamma moments
give
\[
 \E|h|\leq\beta F(\lambda),\qquad
 F(\lambda)=\E\sqrt{\lambda K_3(T)+(1-\lambda)K_1(T)},
\]
where
\[
 K_3(t)=\frac{t^3}{60},\qquad
 K_1(t)=\frac{t(t-4)^2}{12},\qquad
 \E K_3(T)=\E K_1(T)=1.
\]
It remains to show $F(\lambda)<83/100$ uniformly.

Partition the positive half-line into $I_j=[2j,2j+2)$ for $0\leq j<12$
and $I_{12}=[24,\infty)$. Write
\[
 p_j=\Pr(T\in I_j),\quad
 a_j=\E[K_3(T)\mathbf1_{I_j}],\quad
 b_j=\E[K_1(T)\mathbf1_{I_j}].
\]
Cellwise Cauchy--Schwarz yields
\[
 F(\lambda)\leq\sum_{j=0}^{12}
    \sqrt{p_j(\lambda a_j+(1-\lambda)b_j)}.
\]
For any positive rational $\tau_j$, the inequality
$\sqrt z\leq(z/\tau_j+\tau_j)/2$ makes the right-hand side at most
\[
 \frac12\sum_j\tau_j+
 \frac12\max\left\{\sum_j\frac{p_ja_j}{\tau_j},
                       \sum_j\frac{p_jb_j}{\tau_j}\right\}.
\]
The certificate uses the following $10^{12}\tau_j$, in order:
\[
\begin{gathered}
209348125670,\ 269171624541,\ 185197047213,\ 107072713153,\\
41670053327,\ 12718770633,\ 3329218641,\ 784596676,\\
171348368,\ 35321389,\ 6958708,\ 1321845,\ 297387.
\end{gathered}
\]
The survival integrals defining the cells are elementary:
\begin{align*}
 \Pr(T\geq t)&=e^{-t}(1+t+t^2/2),\\
 \E[K_3(T)\mathbf1_{T\geq t}]
 &=e^{-t}\sum_{j=0}^{5}\frac{t^j}{j!},\\
 \E[K_1(T)\mathbf1_{T\geq t}]
 &=\frac{e^{-t}}{24}
 (t^5-3t^4+4t^3+12t^2+24t+24).
\end{align*}
All three multiplying polynomials are positive for $t\geq0$.
At the endpoints $t=2j$ the exponential is $q^j$, where $q=e^{-2}$.
Set
\[
 L=\sum_{j=0}^{40}\frac{2^j}{j!},\qquad
 U=L+\frac{2^{41}}{41!}\frac1{1-2/42}.
\]
The positive exponential series gives $L<e^2<U$, hence $1/U<q<1/L$.
For each cell, use the upper survival bound at its left endpoint and the
lower survival bound at its right endpoint to get rational upper bounds
$\bar p_j,\bar a_j,\bar b_j$. These bounds and the displayed $\tau_j$ are
checked by \texttt{six\_dimensional\_certificate.py}, using only integer
and rational arithmetic, to satisfy
\[
 \frac12\left(\sum_j\tau_j+\sum_j\frac{\bar p_j\bar a_j}{\tau_j}\right)
 <\frac{83}{100},\qquad
 \frac12\left(\sum_j\tau_j+\sum_j\frac{\bar p_j\bar b_j}{\tau_j}\right)
 <\frac{83}{100}.
\]
This proves $F(\lambda)<83/100$ for every $\lambda\in[0,1]$.
Multiplication by $\beta$ gives the asserted constant.
Finally $83\sqrt5/(75\pi)<\sqrt{2/\pi}$ is equivalent to
$\pi>6889/2250$, which follows already from $\pi>223/71$.
\end{proof}
\end{document}
